% One-time imported draft. Editable native source; no Markdown regeneration.
\section{Copy state without copying its representation
mistakes}\label{008-life-generations-p03-copy-state-without-copying-its-representation-mistakes}

The seed is a compact 64-byte input, while a working board has padded
rows. Copying one long block from seed to board base would put states
into halo positions and give the interior the wrong row alignment. The
copy must walk logical rows and skip the two physical halo columns
between successive interior rows.

The source cursor advances eight bytes per logical row. The
working-board cursor also advances through eight interior bytes, then
skips two. After the final row, a cursor can point beyond the final
interior cell without being dereferenced again. The row counter controls
whether another load occurs.

Only \texttt{board\_a} is initialized from the seed. The other interior
may contain old or arbitrary bytes. Computing a new generation must
overwrite every destination cell with a binary result, including dead
states. Assuming the destination begins cleared would make later buffer
reuse depend on stale contents.
